\documentclass[12pt]{article}
\usepackage{ceai}

\begin{document}

\title{CE 488/5588 Homework 06\\
\vspace{2in}
\pmb{Name: \rule{2in}{0.15mm}}
\author{}
}
\maketitle
\newpage

\section*{Topic 10: Generalization and Model Selection}

\begin{ch-box}
This homework is based on the lecture content on generalization, re-sampling methods,
model selection, and ensemble learning. Answer all five questions. Show simple calculations
where requested.
\end{ch-box}

\section*{Conceptual Questions and Short Calculations}

\subsection*{Question 1: Fill in the blanks}
For each statement below, fill in the most appropriate term from the list:

\begin{center}
\textbf{underfitting, overfitting, bias, variance}
\end{center}

\begin{enumerate}[label=(\alph*)]
    \item A model that performs very well on the training set but poorly on unseen data is showing \underline{\hspace{3cm}}.
    \item A model that is too simple to capture the main pattern in the data is showing \underline{\hspace{3cm}}.
    \item Error caused by oversimplifying the true relationship is called \underline{\hspace{3cm}}.
    \item Error caused by excessive sensitivity to the sampled training data is called \underline{\hspace{3cm}}.
\end{enumerate}

\subsection*{Question 2: Multiple choice}
A civil engineer has only 120 labeled inspection cases for a binary classification problem. She wants a more reliable estimate of unseen-data performance than one random train-test split can provide. Which choice below is the \textbf{best} first step?

\begin{itemize}
    \item [A.] Increase the model complexity until the training accuracy reaches 100\%.
    \item [B.] Use K-fold cross-validation on the available labeled data.
    \item [C.] Evaluate performance only on the training set because all data are valuable.
    \item [D.] Lower the classification threshold to 0.1 before any evaluation.
\end{itemize}

Best answer: \rule{1in}{0.15mm}

\bigskip

Briefly explain why your choice is the best one.\\


\vspace{20pt}
\fullunderline

\vspace{20pt}
\fullunderline

\subsection*{Question 3: Cross-validation calculation}
An engineer evaluates one classifier using 5-fold cross-validation. The validation accuracies from the five folds are:

\[
0.81,\quad 0.78,\quad 0.84,\quad 0.80,\quad 0.77
\]

Answer the following:

\begin{enumerate}[label=(\alph*)]
    \item Compute the mean cross-validation accuracy.\\
    \rule{2in}{0.15mm}
    \item Compute the range of the fold accuracies (maximum minus minimum).\\
    \rule{2in}{0.15mm}
    \item Based on these five scores, does the model appear extremely unstable across folds? Answer yes or no and give one short reason.\\
    \fullunderline
\end{enumerate}

\subsection*{Question 4: Threshold selection}
A binary classifier outputs the following class-membership probabilities for four new data points:

\[
0.18,\quad 0.42,\quad 0.67,\quad 0.91
\]

Assume the prediction rule is
\[
\text{predict class }1 \text{ if } P(c\mid \bm{x}) \ge \tau,
\]
where \(\tau\) is the threshold.

\begin{enumerate}[label=(\alph*)]
    \item If \(\tau = 0.50\), write the predicted labels for the four data points.\\
    \rule{3in}{0.15mm}
    \item If \(\tau = 0.30\), write the predicted labels for the four data points.\\
    \rule{3in}{0.15mm}
    \item Which threshold, 0.50 or 0.30, is more likely to increase recall? \underline{\hspace{1.5cm}}.
\end{enumerate}

\subsection*{Question 5: Ensemble methods --- identify the strategy}
Match each scenario below with the most appropriate ensemble term:

\begin{center}
\textbf{bagging, boosting, stacking}
\end{center}

\begin{enumerate}[label=(\alph*)]
    \item Several models are trained on different sampled subsets of the training data, and their predictions are averaged. This is \underline{\hspace{3cm}}.
    \item Weak learners are trained sequentially, and each new learner focuses more on the previous mistakes. This is \underline{\hspace{3cm}}.
    \item A logistic regression model is trained to combine the predictions from an SVM, a decision tree, and a neural network. This is \underline{\hspace{3cm}}.
\end{enumerate}

\bigskip

In one or two sentences, explain why ensemble learning can sometimes improve generalization.\\

\vspace{20pt}
\fullunderline

\vspace{20pt}
\fullunderline

\end{document}
